\begin{longtable}[c]{|l|m{11cm}|}
    \caption{写内存指令}
    \label{tab:pie-store-instructions}\\
    \hline
    \rowcolor{lightgray}
        \textbf{指令} & \textbf{定义} \\\hline
    \endfirsthead
    \rowcolor{lightgray}
    \textbf{指令} & \textbf{定义} \\\hline
    \endhead
    \endfoot
    
    \hyperref[instruction:STQR]{ST.QR} & 将 QR 寄存器中的 16-字节数据写入内存中。
    \\\hline
    \hyperref[instruction:EEVST128XP]{EE.VST.128.XP} & 将 16-字节数据写入内存，随后访存地址加上 AR 寄存器内数值。
    \\\hline
    \hyperref[instruction:EEVST128IP]{EE.VST.128.IP} & 将 16-字节数据写入内存，随后访存地址加上立即数。
    \\\hline
    \hyperref[instruction:EEVSTH64XP]{EE.VST.[H/L].64.XP} & 将 8-字节数据写入内存，随后访存地址加上 AR 寄存器内数值。
    \\\hline
    \hyperref[instruction:EEVSTH64IP]{EE.VST.[H/L].64.IP} & 将 8-字节数据写入内存，随后访存地址加上立即数。
    \\\hline
    \hyperref[instruction:EESTF64XP]{EE.STF.[64/128].XP} & 将 8-字节/16-字节数据写入内存，随后访存地址加上 AR 寄存器内数值。
    \\\hline
    \hyperref[instruction:EESTF64IP]{EE.STF.[64/128].IP} & 将 8-字节/16-字节数据写入内存，随后访存地址加上立即数。
    \\\hline 
    \hyperref[instruction:EESTQACCHH32IP]{EE.ST.QACC\_[H/L].H.32.IP} & 将 4-字节数据写入内存，随后访存地址加上立即数。
    \\\hline
    \hyperref[instruction:EESTQACCHL128IP]{EE.ST.QACC\_[H/L].L.128.IP} & 将 16-字节数据写入内存，随后访存地址加上立即数。
    \\\hline
    \hyperref[instruction:EESTACCXIP]{EE.ST.ACCX.IP} & 将 ACCX 寄存器内数值 0 扩展为 8-字节数据后写入内存，随后访存地址加上立即数。
    \\\hline
    \hyperref[instruction:EESTUASTATEIP]{EE.ST.UA\_STATE.IP} & 将 16-字节数据写入内存，随后访存地址加上立即数。
    \\\hline
    \hyperref[instruction:EESTXQ32]{EE.STXQ.32} & 先更新访存地址，随后将 4-字节数据写入内存。
    \\\hline
\end{longtable}